% Source: https://artofproblemsolving.com/community/c3h2580839_doing_roots_by_hand_replicate_of_my_blog_post
% Stable identity: aops-topic-2580839-q01
For a real number $r$ and a positive integer $k$, define $\dbinom r0=1$ and
\[\dbinom rk=\dfrac{r(r-1)(r-2)\cdots(r-k+1)}{k!}.\]
For $|x|<1$ the binomial series gives $(1+x)^r=\displaystyle\sum_{k\ge0}\dbinom rk x^k$.
\begin{enumerate}
\item[(a)] Compute $\dbinom{\dfrac13}{k}$ for $k=0,1,2,3,4$.
\item[(b)] Use the first five terms of the series to approximate $\sqrt[3]{1.8}$, to four decimal places.
\item[(c)] Explain why the series cannot be used directly for $\sqrt[3]{10}=(1+9)^{\dfrac13}$. Then use $10=8\left(1+\dfrac14\right)$ and the first three terms of the series to approximate $\sqrt[3]{10}$.
\end{enumerate}
